% Candidate subsection generated from the audited 2026-08-27 long run.
% The data pass production audits; the fluctuation-theorem conclusion remains
% provisional and is explicitly reported as not verified in the sampled window.
% Keep separate from the mentor-synchronized draft until the scientific scope
% and placement are approved.

\subsection{Finite-time current fluctuations and thermodynamic diagnostics}
\label{sec:finite-time-fluctuations}

We next examine the finite-time fluctuations of the bulk action current and
their relation to thermodynamic transport.  The simulation is performed in
projection-free Cartesian variables, $c_j=x_j+\mathrm{i}y_j$, using the same
boundary temperatures $T_L=10$ and $T_R=2$ as in the transport study.  After a
burn-in time $B=500$, we record $1{,}000{,}064$ non-overlapping blocks of length
$t=20$ for each $n=10,20,30,40$.  The blocks are generated from 128 independent
streams and are aggregated to averaging windows $t=20,40,\ldots,200$.  In
addition to the bulk action-current average $J_M(t)$, each block stores the
heat $Q_L$ and $Q_R$ delivered by the two bath substeps and the system-energy
change $\Delta E$.  With heat positive into the chain, the medium entropy is
\begin{equation}
  \Sigma_t^{\mathrm m}=-\frac{Q_L}{T_L}-\frac{Q_R}{T_R}.
\end{equation}
The discrete first-law residual $Q_L+Q_R-\Delta E$ is below $10^{-5}$ in RMS
rate for every chain length, all midpoint solves converge, and matched
timestep-halving, equal-temperature, and bath-swapped controls pass.

The production means of the action current are $0.39809$, $0.11979$,
$0.05426$, and $0.03052$ for $n=10,20,30,40$, respectively.  A log--log fit
gives
\begin{equation}
  \langle J_M(n)\rangle = 28.889\,n^{-1.8485},
  \qquad R^2=0.9987,
\end{equation}
in independent agreement with the canonical current calculation.

For a threshold $A$, both empirical action-current tails show an approximately
linear dependence of log probability on the averaging time wherever raw
two-tail resolution is available:
\begin{equation}
  \log \Pr[J_M(t)\ge A]=c_+(A)-t I_+(A),\qquad
  \log \Pr[J_M(t)\le-A]=c_-(A)-t I_-(A).
  \label{eq:finite-time-rate-proxy}
\end{equation}
We require at least 20 observed tail events, probability at most $0.2$, and at
least four resolved averaging times.  All qualified fits have $R^2\ge0.98$,
and the fitted rate proxies $I_\pm(A)$ increase with threshold magnitude.  The
resolved positive-tail intervals are $A\in[0.46,0.81]$, $[0.15,0.33]$,
$[0.08,0.20]$, and $[0.05,0.15]$ for $n=10,20,30,40$, respectively.  The
negative tail is resolved over $A\in[0.01,0.06]$ for $n=20$ and
$A\in[0.01,0.08]$ for $n=30,40$; for $n=10$ it becomes too rare to support a
four-time-point direct-sampling fit.  These results support a finite-time
large-deviation description but do not by themselves imply a fluctuation
theorem.

The central $1$--$99\%$ action-current distribution becomes narrower with
increasing $t$ and is qualitatively Gaussian-like.  Its far tails, however,
are not captured by the Gaussian fixed by the full-sample mean and variance.
At $t=20$, a joint fit of one Gaussian to both empirical tails requires a width
approximately $13\%$, $20\%$, $26\%$, and $28\%$ larger than the corresponding
full-sample width for $n=10,20,30,40$.  Gaussianity is therefore used only as a
shape benchmark.

To test thermodynamic symmetry, define
\begin{equation}
  R_t(a)=\frac{1}{t}\log\frac{p_t(a)}{p_t(-a)},
  \qquad a=\frac{\Sigma_t^{\mathrm m}}{t}.
\end{equation}
The simple asymptotic reference is $R_t(a)=a$.  The directly fitted
medium-entropy slopes remain below unity in all raw-count-qualified windows.
They generally increase with $t$ while both tails remain resolved---for
example, the $n=30$ values are $0.247$, $0.316$, and $0.393$ for
$t=20,40,60$---but do not reach unity before negative events become
resolution-limited.  The bath-heat symmetry slopes likewise remain below the
reference affinity $\Delta\beta=T_R^{-1}-T_L^{-1}=0.4$.  A second symmetric,
equal-width fit with an independently defined quantile range gives the same
qualitative conclusion; the action-current slopes themselves show appreciable
fit-window dependence and are not assigned a thermodynamic target.

This finite-time result must be interpreted with care.  An exact detailed
fluctuation theorem concerns the total trajectory entropy
\begin{equation}
  \Delta s_{\rm tot}=\Sigma_t^{\mathrm m}+\Delta s_{\rm sys},
\end{equation}
where the NESS system-entropy endpoint term is not reconstructed in the
present calculation.  We therefore conclude that the standard thermodynamic
fluctuation symmetry is not verified in the directly sampled window, rather
than interpreting the finite-time slope deficit as a violation.

Finally, heat and action currents become more correlated with increasing
averaging time, but they are not tightly coupled.  At $t=200$, the Pearson
correlation decreases from $0.934$ for $n=10$ to $0.683$ for $n=40$, while the
heat variance unexplained by a linear action-current regression increases from
$0.128$ to $0.533$.  Consequently, the action current cannot be substituted
for thermodynamic entropy production, particularly in the longer chains.
